\DocumentMetadata{
  pdfversion=2.0,
  pdfstandard=ua-2,
  lang=en-US,
  tagging=on,
  tagging-setup={math/setup=mathml-SE,tabsorder=structure}
}
\documentclass[11pt,letterpaper]{article}
\usepackage[margin=0.68in,includefoot]{geometry}
\usepackage{newcomputermodern}
\usepackage{amsmath,amscd,mathtools}
\providecommand{\square}{\mdlgwhtsquare}
\usepackage[hidelinks]{hyperref}
\hypersetup{
  pdftitle={Probability PhD Exam, August 2020},
  pdfauthor={University of Florida Department of Mathematics},
  pdfsubject={Graduate mathematics examination},
  pdfdisplaydoctitle=true,
  bookmarksopen=true
}
\setcounter{secnumdepth}{0}
\setlength{\parindent}{0pt}
\setlength{\parskip}{0.28em}
\setlength{\emergencystretch}{3em}
\setlength{\leftmargini}{2.15em}
\setlength{\leftmarginii}{2.65em}
\setlength{\leftmarginiii}{3em}
\renewcommand{\labelenumi}{\textbf{\theenumi.}}
\pagestyle{empty}
\raggedbottom
\allowdisplaybreaks
\newenvironment{examproblems}[1]{%
  \begin{enumerate}%
  \setcounter{enumi}{#1}%
  \setlength{\topsep}{0.35em}%
  \setlength{\partopsep}{0pt}%
  \setlength{\itemsep}{0.48em}%
  \setlength{\parsep}{0.18em}%
}{\end{enumerate}}
\newenvironment{examsubparts}{%
  \begin{enumerate}%
  \setlength{\topsep}{0.18em}%
  \setlength{\partopsep}{0pt}%
  \setlength{\itemsep}{0.16em}%
  \setlength{\parsep}{0.1em}%
}{\end{enumerate}}
\newenvironment{examinstructions}{%
  \par\begingroup\itshape
}{%
  \par\endgroup\vspace{0.18em}
}
\makeatletter
\renewcommand\section{\@startsection{section}{1}{0pt}%
  {-1.25ex plus -.25ex minus -.1ex}%
  {0.55ex plus .1ex}%
  {\normalfont\large\bfseries}}
\renewcommand{\@maketitle}{%
  \newpage\null\vskip -1.2em%
  {\footnotesize\bfseries
    \href{https://www.ufl.edu/}{University of Florida}, \href{https://math.ufl.edu/}{Department of Mathematics}\par}%
  \vskip 0.18em%
  \tagstructbegin{tag=Title}%
    {\LARGE\bfseries\href{https://gma.math.ufl.edu/exams/probability-phd/}{Probability PhD Exam}, August 2020\par}%
  \tagstructend%
  \vskip 0.35em%
  \tagmcbegin{artifact}%
    \hrule height 1pt%
  \tagmcend%
  \vskip 0.55em%
}
\makeatother
\title{Probability PhD Exam, August 2020}
\author{}
\date{}
\begin{document}
\maketitle
\thispagestyle{empty}
\begin{examinstructions}
Carefully justify your answers.
\end{examinstructions}
\begin{examproblems}{0}
\item
This problem has three parts.
\begin{examsubparts}
\item[\textnormal{1.}]
State the weak and strong law of large numbers for i.i.d. random variables.
\item[\textnormal{2.}]
Prove the weak law of large numbers (for i.i.d. random variables).
\item[\textnormal{3.}]
Prove the strong law of large numbers under the additional assumption of finite fourth moment \(\bigl(\mathbb{E}[X_1^4]<+\infty\bigr)\).
\end{examsubparts}
\item
Let \(X,Y\) be independent and uniformly distributed on \([0,1]\).
\begin{examsubparts}
\item[\textnormal{1.}]
Find the density of the random variable \(X+Y\).
\item[\textnormal{2.}]
Find the density of the random variable \(\mathbb{E}[X\mid X+Y]\).
\end{examsubparts}
\item
Let \(\{X_n\}\) be a sequence of random variables uniformly bounded, that is, there exists \(M>0\) such that for all \(n\in\mathbb{N}\), \(|X_n|\le M\). Prove that \(\{X_n\}\) converges to~\(0\) in \(L^1\) if and only if \(\{X_n\}\) converges to~\(0\) in probability.
\item
Let \(\{X_n\}\) be a sequence of random variables such that for all \(n\ge 1\), \(X_n\) has a Poisson distribution of parameter~\(n\). Does the random variable 
\[
\frac{X_n-n}{\sqrt n}
\]
 converges in distribution as \(n\to+\infty\)? If so, what is the limit? What about convergence in probability? Carefully justify all your answers.
\item
Let \(\{X_n\}\) be a sequence of i.i.d. random variables in \(L^1\) \(\bigl(\mathbb{E}[|X_1|]<+\infty\bigr)\). Take \(\mathcal{F}_n=\sigma\{X_1,\ldots,X_n\}\), \(n\ge 1\), and denote \(S_n=\sum_{k=1}^n X_k\), \(S_0=0\).
\begin{examsubparts}
\item[\textnormal{1.}]
Let~\(\tau\) be an \(\{\mathcal{F}_n\}\)-stopping time in \(L^1\). Prove that \(S_\tau\) is in \(L^1\), and 
\[
\mathbb{E}[S_\tau]=\mathbb{E}[X_1]\mathbb{E}[\tau].
\]
\item[\textnormal{2.}]
Let \(\{X_n\}\) be a sequence of i.i.d. random variables such that \(\mathbb{P}(X_1=1)=\mathbb{P}(X_1=-1)=\frac12\). Denote \(S_n=\sum_{k=1}^n X_k\), \(S_0=0\), and \(T=\inf\{n\ge 0:S_n=1\}\). Prove that \(\mathbb{E}[T]=+\infty\). (One may use question 1.)
\end{examsubparts}
\item
Let \(X_1,\ldots,X_n\) be a sequence of i.i.d. random variables in \(L^1\). Compute \(\mathbb{E}[X_1\mid X_1+\cdots+X_n]\).

\par
Hint: for all \(i,j\in\{1,\ldots,n\}\), \(\mathbb{E}[X_i\mid X_1+\cdots+X_n]=\mathbb{E}[X_j\mid X_1+\cdots+X_n]\).
\item
This problem has two parts.
\begin{examsubparts}
\item[\textnormal{1.}]
Give the definition of a standard Brownian motion.
\item[\textnormal{2.}]
Let \(\{B_t\}\) and \(\{\widetilde B_t\}\) be two independent standard Brownian motion. Let \(\rho\in(0,1)\). Define, for \(t\ge 0\), 
\[
X_t=\rho B_t+\sqrt{1-\rho^2}\,\widetilde B_t.
\]
 Is \(\{X_t\}\) a standard Brownian motion? Carefully justify.
\end{examsubparts}
\item
Let \(\{B_t\}\) be a standard Brownian motion.
\begin{examsubparts}
\item[\textnormal{1.}]
Fix \(u>0\). Prove that \(M_t=e^{uB_t-\frac12u^2t}\) is a martingale (with respect to the same filtration as \(\{B_t\}\)).
\item[\textnormal{2.}]
Fix \(a>0\). Define \(T_a=\inf\{t\ge 0:B_t=a\}\). Prove that \(T_a<+\infty\) a.s. and \(\mathbb{E}[T_a]=+\infty\).

\par
Hint: One may find the density of \(T_a\) by considering \(X_t=\max_{0\le s\le t}B_s\), and noticing that \(\{T_a\le t\}=\{X_t\ge a\}\).
\end{examsubparts}
\end{examproblems}
\end{document}
